% ID: SEL-D0066-OB-E04
% ESTADO: cotejada. Segunda pasada visual y borrador compilado revisados.
% Fuente: Selectividad/sel_2010_matematicas_aplicadas/mat_apli/mat_apli_sept/mat_apli_sept_rb/EXAMEN-ANDALUCÍA-2.pdf
% SHA-256: 24c494a5d013dbc7f5adaac9896dee2d79f22540518359c67eac29945a26d4f0
% Página: 2; opción B; ejercicio original 4.

\textbf{(2.5 puntos)} Se sabe que los años de vida de los individuos de una población es una variable aleatoria Normal con desviación típica $8.9$ años. Una muestra aleatoria de 100 individuos de esa población mostró una vida media de $71.8$ años. Mediante un contraste de hipótesis unilateral, \textquestiondown{}puede afirmarse con los datos anteriores que la vida media es mayor de 70 años, a un nivel de significación $\alpha=0.05$?
